\chapter{唯一的快速输出}
% 完整版：chapters/13_muscles.tex

你在世界上快速做出的一切——一句话、一个眼神、一步——都是肌肉的收缩。机器还有第二个输出，缓慢的、化学的，那是第十六章的内容。这一章讲肌肉。

链条的最后一环是\nt{ACET}神经元。每根肌纤维只听一个这样的神经元，而一个神经元控制一组纤维，从几百根到两千来根。有意思的是，神经与肌肉的连接是全身最可靠的连接。在大脑内部，连接会丢失一半以上的信号，而在这里，每一声咔嗒释放的物质都比需要的多好几倍。在输出端，一次错误的代价是一次动作，所以可靠性是用充足的余量换来的。

\section*{鼓点变成推力}

一声咔嗒——肌肉一次短暂的抽动。咔嗒稀疏时，肌肉一下一下地抽动。咔嗒密集时，抽动就融合成平稳的力量，正如密集的鼓点融合成一片轰鸣。肌肉把脉冲流变成力，就像一个最简单的数模转换器。

\pencilfig{113}{%
% twitch = alpha function; the sum of twitches for spikes 1 s and 0.16 s apart (arbitrary units of time)
\begin{scope}[x=0.95cm, y=0.36cm]
\foreach \dx/\t in {0/稀疏,4.6/密集}{
  \begin{scope}[xshift=\dx cm]
  \draw[pen] (0,0) -- (4,0); \node[font=\small\itshape, text=pcGray, anchor=south] at (2,5.4) {\t};
  \end{scope}}
\draw[penlite=pcRed] plot[smooth] coordinates {(0.00,0.00) (0.05,0.00) (0.10,0.00) (0.15,0.00) (0.20,0.00) (0.25,0.45) (0.30,0.73) (0.35,0.90) (0.40,0.98) (0.45,1.00) (0.50,0.98) (0.55,0.94) (0.60,0.88) (0.65,0.81) (0.70,0.74) (0.75,0.66) (0.80,0.59) (0.85,0.52) (0.90,0.46) (0.95,0.41) (1.00,0.35) (1.05,0.31) (1.10,0.27) (1.15,0.23) (1.20,0.20) (1.25,0.62) (1.30,0.88) (1.35,1.02) (1.40,1.08) (1.45,1.09) (1.50,1.06) (1.55,1.00) (1.60,0.93) (1.65,0.86) (1.70,0.78) (1.75,0.70) (1.80,0.62) (1.85,0.55) (1.90,0.48) (1.95,0.42) (2.00,0.37) (2.05,0.32) (2.10,0.28) (2.15,0.24) (2.20,0.21) (2.25,0.62) (2.30,0.88) (2.35,1.03) (2.40,1.09) (2.45,1.09) (2.50,1.06) (2.55,1.01) (2.60,0.94) (2.65,0.86) (2.70,0.78) (2.75,0.70) (2.80,0.62) (2.85,0.55) (2.90,0.48) (2.95,0.42) (3.00,0.37) (3.05,0.32) (3.10,0.28) (3.15,0.24) (3.20,0.21) (3.25,0.62) (3.30,0.88) (3.35,1.03) (3.40,1.09) (3.45,1.09) (3.50,1.06) (3.55,1.01) (3.60,0.94) (3.65,0.86) (3.70,0.78) (3.75,0.70) (3.80,0.62) (3.85,0.55) (3.90,0.48) (3.95,0.42) (4.00,0.37)};
\foreach \s in {0.20,1.20,2.20,3.20}{\draw[pcBlue, line width=0.8pt] (\s,-0.35) -- (\s,-1.2);}
\begin{scope}[xshift=4.6cm]
\draw[penlite=pcRed] plot[smooth] coordinates {(0.00,0.00) (0.05,0.00) (0.10,0.00) (0.15,0.00) (0.20,0.00) (0.25,0.45) (0.30,0.73) (0.35,0.90) (0.40,1.35) (0.45,1.68) (0.50,1.85) (0.55,2.19) (0.60,2.51) (0.65,2.64) (0.70,2.84) (0.75,3.13) (0.80,3.22) (0.85,3.27) (0.90,3.55) (0.95,3.62) (1.00,3.54) (1.05,3.82) (1.10,3.88) (1.15,3.80) (1.20,3.98) (1.25,4.06) (1.30,3.97) (1.35,4.07) (1.40,4.16) (1.45,4.09) (1.50,4.10) (1.55,4.23) (1.60,4.17) (1.65,4.10) (1.70,4.26) (1.75,4.23) (1.80,4.07) (1.85,4.27) (1.90,4.27) (1.95,4.12) (2.00,4.26) (2.05,4.29) (2.10,4.17) (2.15,4.24) (2.20,4.31) (2.25,4.21) (2.30,4.21) (2.35,4.31) (2.40,4.24) (2.45,4.16) (2.50,4.31) (2.55,4.27) (2.60,4.10) (2.65,4.30) (2.70,4.29) (2.75,4.15) (2.80,4.28) (2.85,4.31) (2.90,4.18) (2.95,4.25) (3.00,4.32) (3.05,4.22) (3.10,4.21) (3.15,4.32) (3.20,4.25) (3.25,4.17) (3.30,4.31) (3.35,4.27) (3.40,4.11) (3.45,3.86) (3.50,3.56) (3.55,3.25) (3.60,2.93) (3.65,2.63) (3.70,2.33) (3.75,2.06) (3.80,1.81) (3.85,1.58) (3.90,1.38) (3.95,1.19) (4.00,1.03)};
\foreach \s in {0.20,0.36,0.52,0.68,0.84,1.00,1.16,1.32,1.48,1.64,1.80,1.96,2.12,2.28,2.44,2.60,2.76,2.92,3.08,3.24}{\draw[pcBlue, line width=0.8pt] (\s,-0.35) -- (\s,-1.2);}
\end{scope}
\node[lbl, text=pcRed, anchor=west] at (1.2,1.9) {抽动};
\node[lbl, text=pcRed, anchor=south] at (6.6,4.55) {平稳的力};
\node[lbl, text=pcBlue, anchor=north] at (4.3,-1.35) {神经元的咔嗒};
\end{scope}
}{稀疏的咔嗒产生一次次单独的肌肉抽动；密集的咔嗒融合成平稳而大得多的力。}

\section*{浮点数表示的力}

大脑有两个控制力量的旋钮：神经元咔嗒得多频繁，以及有多少组纤维被启用。各组按顺序启用——从最弱到最强，而这个顺序不需要谁来计算：小神经元只是更容易兴奋。结果，每新加入一组，都会在已有的力上增加大致相同的\emph{比例}。拿鸡蛋时，力是按极小的份量来调配的；提手提箱时，则按大份量。工程师会从中认出浮点数：数量级由启用的组数决定，精确值则由咔嗒频率决定。代价是：动作越用力，就越不精确。按照一个很有影响的假说，这就是我们的动作如此平滑的原因：平滑的指令产生的偏差最小。

\section*{只能拉，不能推}

肌肉只会拉。所以每个关节由一对朝相反方向拉的肌肉负责——又一次用两根导线代替符号。两者力量之差使关节转动，力量之和使关节更僵硬。端着一杯满满的茶时，你同时绷紧两块肌肉：转动不变，但手臂更能扛住冲撞。

\pencilfig{13}{\pencilart{13}}{肌肉只会拉，所以要两块：力之差让关节转动，力之和让关节更硬。}

\section*{迟到的反馈}

肌肉里有拉伸传感器，最短的回路直接在脊髓里闭合：肌肉被拉长——它自己就收缩，与此同时对抗肌被关掉。关掉它的是\nt{GLYC}类型的抑制性神经元：这种类型总算找到了用武之地。

但每条回路都有延迟——在脊髓里是几十毫秒，经由眼睛是几百毫秒。而迟到的反馈会让系统摇摆起来，就像一个看着滞后的路况打方向盘的司机。回路越长，纠正就必须越慢。脊髓回路开始振荡的界限大约是每秒八次。这接近每个健康人都有的细微颤抖的频率，而拉伸回路是它可能的来源之一。

那我们怎么能动得又快又准？按照一种流行的假说——靠预测：大脑保存着一个身体的内部模型，不等传感器回报就算出指令的结果。

\section*{没有节拍器的节律}

走路和呼吸都是有节律的，却不需要节拍器。两群互相抑制、又会逐渐疲劳的神经元，会自己轮流工作：胜者累了，休息过的对手占了上风，如此循环往复。来自上方的一条恒定指令——“走”——就地变成了“左、右、左、右”的节律。

\fullversion{13}
